% Generated draft: requires mathematical and source review.
\documentclass[11pt]{article}
\usepackage[margin=25mm]{geometry}
\usepackage{fontspec}
\setmainfont{Latin Modern Roman}
\usepackage{amsmath,amssymb,mathtools,unicode-math}
\setmathfont{Latin Modern Math}
\usepackage{enumitem,needspace,xurl,hyperref}
\setlength{\parindent}{0pt}
\setlength{\parskip}{5pt}
\setlength{\emergencystretch}{3em}
\newcommand{\N}{\mathbb{N}}
\newcommand{\Z}{\mathbb{Z}}
\newcommand{\Q}{\mathbb{Q}}
\newcommand{\R}{\mathbb{R}}
\newcommand{\C}{\mathbb{C}}
\newcommand{\F}{\mathbb{F}}
\newcommand{\A}{\mathbb{A}}
\newcommand{\PP}{\mathbb{P}}
\newcommand{\E}{\mathbb{E}}
\newcommand{\eps}{\varepsilon}
\newcommand{\dd}{\,\mathrm{d}}
\begin{document}
\section*{UnsolvedMath ID: 30003597}
\textbf{Source problem code:} catalogue entry
\subsection{Definitions and mathematical statement}
Let \(d\in\Z\) with \(d\geq4\), let \(N\) be a lattice of rank \(d\), and put \(N_{\R}=N\otimes_{\Z}\R\). A \(d\)-dimensional lattice polytope is the convex hull of finitely many points of \(N\), with affine span \(N_{\R}\). Such a polytope \(P\) is called canonical Fano if
\[
\operatorname{int}(P)\cap N=\{0\}.
\]
Use lattice-relative ordinary volume \(\operatorname{vol}_N\), for which a fundamental parallelepiped of \(N\) has volume \(1\). The supplied record does not explicitly specify its volume normalization; this convention, consistent with the displayed denominator \(d!\), requires confirmation against the source.

Define the Sylvester sequence by
\[
s_1=2,\qquad s_{i+1}=\prod_{j=1}^{i}s_j+1\quad(i\geq1).
\]
The proposed assertion is that every such \(P\) satisfies
\[
\operatorname{vol}_N(P)\leq\frac{2(s_d-1)^2}{d!},
\]
with equality only for the canonical Fano simplex denoted \(R_{(d)}\) in the supplied statement (written \(R(d)\) in its original-statement field).

The equality clause remains incompletely specified: the supplied record gives no definition or vertices of \(R_{(d)}\) and does not specify the equivalence relation intended by ``only for.'' Thus the inequality is formulated under the stated volume convention, while a precise identification of its equality cases requires additional source information. The source asks for a sharp bound and equality characterization; these are proposed assertions, not conclusions established here.
\subsection{Short English statement}
For every integer \(d\geq4\), the problem asks whether a \(d\)-dimensional lattice polytope whose unique interior lattice point is the origin has volume at most \(2(s_d-1)^2/d!\), with equality only for a designated simplex \(R_{(d)}\). The record does not define that simplex or the intended equivalence of equality cases, and its volume convention needs confirmation.
\subsection{Sources}
\begin{itemize}
\item Supplied local UnsolvedMath registry record, permanent ID \(30003597\), title \emph{Maximum Volume of Canonical Fano Polytopes}, source problem code \texttt{catalogue entry}, in the Oberwolfach Reports Open Problems collection. No individual proposer is supplied; the proposed year is recorded as 2017.
\item Supplied report citation: \emph{Mini-Workshop: Lattice Polytopes: Methods, Advances, Applications} (2018), Oberwolfach Reports, pp.\ 2659--2701, DOI \texttt{10.4171/owr/2017/44}. Supplied source URL: \url{https://doi.org/10.4171/owr/2017/44}.
\item Supplied literature reference: Balletti--Kasprzyk--Nill, ``maximum dual volume,'' \url{https://doi.org/10.1017/fms.2022.93}.
\item Supplied literature reference: Averkov--Krumpelmann--Nill, ``largest integral simplices,'' \url{https://arxiv.org/abs/1309.7967}.
\item Supplied literature reference: ``original OWR report,'' \url{https://doi.org/10.4171/OWR/2017/44}.
\item Literature assessment dated exactly \texttt{2026-08-21}: the supplied curation labels the problem open and reports that the canonical Fano volume bound is proved for canonical Fano simplices, that the dual-volume analogue is proved in general, and that the sharp primal-volume bound and equality characterization remain open for general canonical Fano polytopes. This is supplied dated curation derived from the OWR collection, not an independently verified literature assessment; its status and cited claims require independent verification.
\end{itemize}
\end{document}
